\documentclass[12pt]{article}

\usepackage[letterpaper,top=2cm,bottom=2cm,left=2.54cm,right=2.54cm,marginparwidth=1.75cm]{geometry}
\usepackage{amsmath}
\usepackage{amssymb}
\usepackage{lineno}
% External equation and appendix references remain clickable in the standalone appendix.
\usepackage[colorlinks=true, allcolors=blue]{hyperref}
\usepackage{comment}
\usepackage{xcolor}
\usepackage{natbib}
\usepackage{xr-hyper}
\externaldocument[A-]{build/appendix_a_cov_formula_dyad}[appendix_a_cov_formula_dyad.pdf]

% The bibliography is printed only in the combined manuscript, so standalone
% appendix citations have labels but no local hyperlink destinations.
\makeatletter
\def\hyper@natlinkstart#1{}
\def\hyper@natlinkend{}
\def\hyper@natlinkbreak#1#2{#1}
\makeatother

\setlength{\parindent}{0pt}
\setlength{\parskip}{0.5em}

% Reduced space above subsection headings (the article default is 3.25ex plus 1ex minus .2ex); format and space below are the defaults.
\makeatletter
\renewcommand\subsection{\@startsection{subsection}{2}{\z@}{-1.5ex plus -.5ex minus -.2ex}{1.5ex plus .2ex}{\normalfont\large\bfseries}}
\makeatother

\newcommand{\vc}{\operatorname{vec}}
\newcommand{\vech}{\operatorname{vech}}
\newcommand{\tr}{\operatorname{tr}}
\newcommand{\R}{\mathbb{R}}
\newcommand{\E}{\mathbb{E}}
\newcommand{\Gap}{\operatorname{Gap}}
\newcommand{\RL}{\mathbf{R}_{\boldsymbol{\Lambda}}}

% This appendix is embedded into the main manuscript via \includepdf, whose
% own consolidated References section is the single printed bibliography
% for the whole document. This file therefore carries no
% \bibliographystyle/\bibliography call. Its \citep/\citet keys resolve
% through bibcites.tex (loaded below).

% Citation labels shared with the manuscript.
\AtBeginDocument{\bibcite{Drysdale2017RestingStateBiotypes}{{1}{2017}{{Drysdale et~al.}}{{}}}
\bibcite{Folland1999}{{2}{1999}{{Folland}}{{}}}
\bibcite{Friston2011}{{3}{2011}{{Friston}}{{}}}
\bibcite{Friston2014}{{4}{2014{a}}{{Friston et~al.}}{{Friston, Kahan, Biswal, and Razi}}}
\bibcite{FristonNodesModes2014}{{5}{2014{b}}{{Friston et~al.}}{{Friston, Kahan, Razi, Stephan, and Sporns}}}
\bibcite{Galioulline2023}{{6}{2023}{{Galioulline et~al.}}{{Galioulline, Fr{\"a}ssle, Harrison, Pereira, Heinzle, and Stephan}}}
\bibcite{Gradziuk2022ColoredNoise}{{7}{2022}{{Gradziuk et~al.}}{{Gradziuk, Torregrosa, and Broedersz}}}
\bibcite{Higham2008}{{8}{2008}{{Higham}}{{}}}
\bibcite{HornJohnson2013}{{9}{2013}{{Horn and Johnson}}{{}}}
\bibcite{Marquardt2007MFICARMA}{{10}{2007}{{Marquardt}}{{}}}
\bibcite{Mityagin2020ZeroSet}{{11}{2020}{{Mityagin}}{{}}}
\bibcite{Novelli2024}{{12}{2024}{{Novelli et~al.}}{{Novelli, Friston, and Razi}}}
\bibcite{Pavliotis2014}{{13}{2014}{{Pavliotis}}{{}}}
\bibcite{Razi2015}{{14}{2015}{{Razi et~al.}}{{Razi, Kahan, Rees, and Friston}}}
\bibcite{SteinShakarchi2003}{{15}{2003}{{Stein and Shakarchi}}{{}}}
}
\begin{document}
\pagestyle{empty}
\linenumbers

\section*{Appendix B: Deriving the Covariance Formula and A Lyapunov Equation for $n\times n$ Networks}
\renewcommand{\theequation}{B-\arabic{equation}}
\setcounter{equation}{0}

% =====================================================================
\subsection*{Setting}
% =====================================================================

Let $\mathbf A$ be a real Hurwitz-stable $n\times n$ matrix and consider
\begin{equation}
    \dot{\boldsymbol x}=\mathbf A\boldsymbol x+\boldsymbol w,
    \qquad
    \mathbf S_w(\omega)=\alpha|\omega|^{-\beta}\boldsymbol{\Sigma}_w,
    \label{eq:sde}
\end{equation}
where $\alpha>0$, $-1<\beta<1$, and the constant spatial noise factor satisfies
$\boldsymbol{\Sigma}_w=\boldsymbol{\Sigma}_w^\top\succeq0$.
The isotropic case is $\boldsymbol{\Sigma}_w=\mathbf I$.
As in \hyperref[A-eq:D_Sigma]{Appendix~A}, a power law over the whole frequency axis is understood as generalized noise;
Gaussianity is not needed for these covariance calculations.
We write $\mathcal U$ for the set of real stable matrices and
$\mathbf R(\omega;\mathbf A)=(i\omega\mathbf I-\mathbf A)^{-1}$ for the resolvent.
With Fourier kernel $e^{-i\omega t}$, the state equation gives
$\hat{\boldsymbol x}=\mathbf R\hat{\boldsymbol w}$, hence
$\mathbf S_x=\mathbf R\mathbf S_w\mathbf R^\dagger$.
The inverse Fourier transform at zero lag, with normalization $1/(2\pi)$, gives the neuronal covariance
\begin{equation}
    \boldsymbol{\Sigma}(\mathbf{A},\beta)=\frac{\alpha}{2\pi}\int_{-\infty}^{\infty}|\omega|^{-\beta}\,
    \mathbf{R}(\omega;\mathbf{A})\,\boldsymbol{\Sigma}_w\,\mathbf{R}(\omega;\mathbf{A})^{\dagger}\,d\omega ,
    \label{eq:Sigma_Sw}
\end{equation}
For $n=2$ and $\boldsymbol{\Sigma}_w=\mathbf I$, this is \eqref{A-eq:D_Sigma}.
Throughout this appendix, $\boldsymbol{\Sigma}$ denotes neuronal covariance.
For a fixed real, frequency-independent observation map $\boldsymbol y=\mathbf H\boldsymbol x$ without observation noise,
the observed covariance is $\boldsymbol{\Sigma}_y=\mathbf H\boldsymbol{\Sigma}\mathbf H^\top$;
this relation requires no invertibility of $\mathbf H$.

% =====================================================================
\subsection*{Reducing the matrix-valued complex integral in \ref{eq:Sigma_Sw} to a series of scalar integrals}
% =====================================================================

Let $\mathbf{A}\in\mathcal{U}$ be diagonalizable, $\mathbf{A}=\mathbf{P}\boldsymbol{\Lambda}\mathbf{P}^{-1}$, where the columns of the
invertible complex matrix $\mathbf{P}$ are eigenvectors of $\mathbf{A}$ and $\boldsymbol{\Lambda}=\operatorname{diag}(\lambda_1,\dots,\lambda_n)$
holds the eigenvalues, each with $\operatorname{Re}\lambda_k<0$. The resolvent inherits this factorization. Since $\mathbf{I}=\mathbf{P}\mathbf{P}^{-1}$, multiplying by the scalar $i\omega$ gives
$i\omega\mathbf{I}=\mathbf{P}\,(i\omega\mathbf{I})\,\mathbf{P}^{-1}$, and subtracting $\mathbf{A}=\mathbf{P}\boldsymbol{\Lambda}\mathbf{P}^{-1}$ from it gives
\begin{equation}
    i\omega\mathbf{I}-\mathbf{A}=\mathbf{P}\,(i\omega\mathbf{I})\,\mathbf{P}^{-1}-\mathbf{P}\boldsymbol{\Lambda}\mathbf{P}^{-1}
    =\mathbf{P}\,(i\omega\mathbf{I}-\boldsymbol{\Lambda})\,\mathbf{P}^{-1}.
    \label{eq:resolvent_similarity}
\end{equation}
The diagonal matrix $i\omega\mathbf{I}-\boldsymbol{\Lambda}$ carries the entries $i\omega-\lambda_k$, and none of them vanishes, because
$\operatorname{Re}(i\omega-\lambda_k)=-\operatorname{Re}\lambda_k>0$ for real $\omega$. It is therefore invertible, as is $\mathbf{P}$. Since the
inverse of a product of invertible matrices is the product of the inverses in reverse order, the resolvent factorizes as
\begin{equation}
    \mathbf{R}(\omega;\mathbf{A})=(i\omega\mathbf{I}-\mathbf{A})^{-1}
    =\mathbf{P}\,(i\omega\mathbf{I}-\boldsymbol{\Lambda})^{-1}\mathbf{P}^{-1}
    =\mathbf{P}\,\RL(\omega)\,\mathbf{P}^{-1}.
    \label{eq:R_eig}
\end{equation}
where $\RL(\omega):=\mathbf{R}(\omega;\boldsymbol{\Lambda})$ is the resolvent of $\boldsymbol{\Lambda}$:
\begin{equation}
    \RL(\omega)=(i\omega\mathbf{I}-\boldsymbol{\Lambda})^{-1}
    =\operatorname{diag}\Bigl(\frac{1}{i\omega-\lambda_1},\dots,\frac{1}{i\omega-\lambda_n}\Bigr).
    \label{eq:RL_eig}
\end{equation}
The diagonal form in \eqref{eq:RL_eig} holds because a diagonal matrix is inverted entry by entry.

Taking the conjugate transpose of \eqref{eq:R_eig} reverses the order of the three factors and applies the operation to each,
\begin{equation}
    \mathbf{R}^{\dagger}=\bigl(\mathbf{P}\,\RL\,\mathbf{P}^{-1}\bigr)^{\dagger}
    =(\mathbf{P}^{-1})^{\dagger}\,\RL^{\dagger}\,\mathbf{P}^{\dagger}
    =\mathbf{P}^{-\dagger}\,\RL^{\dagger}\,\mathbf{P}^{\dagger}.
    \label{eq:Rdag_eig}
\end{equation}
Here $\mathbf{P}^{-\dagger}:=(\mathbf{P}^{-1})^{\dagger}$.
The conjugate transpose of a diagonal matrix is diagonal with its entries conjugated, and since $\omega$ is real the conjugate of
$i\omega-\lambda_l$ is $-i\omega-\bar\lambda_l$, so
\begin{equation}
    \RL^{\dagger}=\operatorname{diag}\Bigl(\frac{1}{-i\omega-\bar\lambda_1},\dots,\frac{1}{-i\omega-\bar\lambda_n}\Bigr).
    \label{eq:RLdag_eig}
\end{equation}
Substituting both factorizations into the integrand of \eqref{eq:Sigma_Sw} gives a product of seven factors, which regroups as
\begin{equation}
    \mathbf{R}\boldsymbol{\Sigma}_w\mathbf{R}^{\dagger}
    =\bigl(\mathbf{P}\,\RL\,\mathbf{P}^{-1}\bigr)\,\boldsymbol{\Sigma}_w\,\bigl(\mathbf{P}^{-\dagger}\,\RL^{\dagger}\,\mathbf{P}^{\dagger}\bigr)
    =\mathbf{P}\,\RL\,\underbrace{\mathbf{P}^{-1}\,\boldsymbol{\Sigma}_w\,\mathbf{P}^{-\dagger}}_{=:\;\mathbf{M}_w}\,\RL^{\dagger}\,\mathbf{P}^{\dagger} .
    \label{eq:RRdag_regroup}
\end{equation}
The three middle factors define the noise spectrum's spatial covariance in eigenvector coordinates.
Even for $\boldsymbol{\Sigma}_w=\mathbf I$, the product $\mathbf P^{-1}\mathbf P^{-\dagger}$ is generally not the identity:
$\mathbf P^{-\dagger}$ is the conjugate transpose of $\mathbf P^{-1}$, not its inverse.
Thus isotropic input in the original coordinates need not be isotropic in eigenvector coordinates. Hence
\begin{equation}
    \mathbf{R}\boldsymbol{\Sigma}_w\mathbf{R}^{\dagger}
    =\mathbf{P}\RL\mathbf{P}^{-1}\boldsymbol{\Sigma}_w\mathbf{P}^{-\dagger}\RL^{\dagger}\mathbf{P}^{\dagger}
    =\mathbf{P}\bigl(\RL\mathbf{M}_w\RL^{\dagger}\bigr)\mathbf{P}^{\dagger}.
    \label{eq:RRdag_eig}
\end{equation}
Here $\mathbf{M}_w:=\mathbf{P}^{-1}\boldsymbol{\Sigma}_w\mathbf{P}^{-\dagger}$. The matrix $\mathbf{M}_w$ does not depend on $\omega$, since $\boldsymbol{\Sigma}_w$, $\mathbf{P}$ and $\mathbf{P}^{-1}$ are all constant.
The frequency therefore enters the integrand only through the two diagonal factors $\RL$ and $\RL^{\dagger}$. To read off the entries of the inner triple product, abbreviate the two diagonal factors as
\begin{equation}
    \RL=\operatorname{diag}(r_1,\dots,r_n),
    \qquad r_k=\frac{1}{i\omega-\lambda_k}.
    \label{eq:RL_entries}
\end{equation}
\begin{equation}
    \RL^{\dagger}=\operatorname{diag}(s_1,\dots,s_n),
    \qquad s_l=\frac{1}{-i\omega-\bar\lambda_l}.
    \label{eq:RLdag_entries}
\end{equation}
A diagonal matrix has the entries $(\RL)_{ka}=r_k\delta_{ka}$ and $(\RL^{\dagger})_{bl}=s_l\delta_{bl}$, where $\delta$ is the Kronecker delta.
The $(k,l)$ entry of the triple product is then the double sum
\begin{equation}
    \begin{aligned}
        \bigl(\RL\mathbf{M}_w\RL^{\dagger}\bigr)_{kl}
        &=\sum_{a=1}^{n}\sum_{b=1}^{n}(\RL)_{ka}\,(M_w)_{ab}\,(\RL^{\dagger})_{bl}\\
        &=\sum_{a=1}^{n}\sum_{b=1}^{n}r_k\,\delta_{ka}\,(M_w)_{ab}\,s_l\,\delta_{bl}
        =r_k\,(M_w)_{kl}\,s_l .
    \end{aligned}
    \label{eq:triple_product_entry}
\end{equation}
because $\delta_{ka}$ is nonzero only when $a=k$ and $\delta_{bl}$ is nonzero only when $b=l$. Multiplying by a diagonal
matrix on the left scales the rows and multiplying by one on the right scales the columns, so each entry of $\mathbf{M}_w$ is
simply rescaled. Collecting the two frequency-dependent factors into
\begin{equation}
    g_{kl}(\omega):=\frac{1}{(i\omega-\lambda_k)(-i\omega-\bar\lambda_l)} ,
    \label{eq:g_def}
\end{equation}
which is the product $r_ks_l$ of the $k$th diagonal entry of $\RL$ and the $l$th diagonal entry of $\RL^{\dagger}$. Thus
\begin{equation}
    \bigl(\RL\mathbf{M}_w\RL^{\dagger}\bigr)_{kl}=(M_w)_{kl}\,g_{kl}(\omega).
    \label{eq:triple_product_g}
\end{equation}

Substituting \eqref{eq:RRdag_eig} into \eqref{eq:Sigma_Sw} gives
\begin{equation}
    \begin{aligned}
    \boldsymbol{\Sigma}(\mathbf{A},\beta)
    &=\frac{\alpha}{2\pi}\int_{-\infty}^{\infty}|\omega|^{-\beta}\,
    \mathbf{P}\bigl(\RL\mathbf{M}_w\RL^{\dagger}\bigr)\mathbf{P}^{\dagger}\,d\omega \\
    &=\frac{\alpha}{2\pi}\,\mathbf{P}
    \Bigl[\int_{-\infty}^{\infty}|\omega|^{-\beta}\,\RL\mathbf{M}_w\RL^{\dagger}\,d\omega\Bigr]\mathbf{P}^{\dagger}.
    \end{aligned}
    \label{eq:Sigma_factor_P}
\end{equation}
A matrix integral is taken entry by entry, and $\mathbf{P}$ and $\mathbf{P}^{\dagger}$ do not depend on $\omega$, so both pass
outside the integral.
Substituting \eqref{eq:triple_product_g} into \eqref{eq:Sigma_factor_P}, the $(k,l)$ entry of the bracket is
\begin{equation}
    \int_{-\infty}^{\infty}|\omega|^{-\beta}\,(M_w)_{kl}\,g_{kl}(\omega)\,d\omega
    =(M_w)_{kl}\int_{-\infty}^{\infty}|\omega|^{-\beta}g_{kl}(\omega)\,d\omega.
    \label{eq:entry_integral}
\end{equation}
since $(M_w)_{kl}$ is a constant that passes through the integral. A matrix whose $(k,l)$ entry is the product of the $(k,l)$
entries of two matrices is their entrywise product, so the bracket term in \eqref{eq:Sigma_factor_P} equals $\mathbf{M}_w\circ\mathbf{Q}$ and
\begin{equation}
\boxed{\begin{aligned}
    \boldsymbol{\Sigma}(\mathbf{A},\beta)
    &=\frac{\alpha}{2\pi}\,\mathbf{P}\,(\mathbf{M}_w\circ\mathbf{Q})\,\mathbf{P}^{\dagger},\\
    q_{kl}
    &:=\int_{-\infty}^{\infty}|\omega|^{-\beta}g_{kl}(\omega)\,d\omega
    =\int_{-\infty}^{\infty}
    \frac{|\omega|^{-\beta}}{(i\omega-\lambda_k)(-i\omega-\bar\lambda_l)}\,d\omega .
\end{aligned}}
    \label{eq:Sigma_hadamard}
\end{equation}
where $\mathbf{Q}$ is the $n\times n$ matrix with entries $q_{kl}$ and $\circ$ denotes the entrywise product.

Each $g_{kl}$ is integrable against $|\omega|^{-\beta}$. Neither factor of its denominator vanishes for real $\omega$, because
$i\omega-\lambda_k=0$ would require $\operatorname{Re}\lambda_k=0$ and $-i\omega-\bar\lambda_l=0$ would require
$\operatorname{Re}\lambda_l=0$, and stability excludes both. The function $g_{kl}$ is therefore continuous on $\R$, and so is
$(1+\omega^2)\,|g_{kl}(\omega)|$. As $|\omega|\to\infty$ the two factors of the denominator satisfy
$|i\omega-\lambda_k|/|\omega|\to1$ and $|{-i\omega}-\bar\lambda_l|/|\omega|\to1$, so $\omega^2|g_{kl}(\omega)|\to1$ and hence
$(1+\omega^2)\,|g_{kl}(\omega)|\to1$. A continuous function on $\R$ with a finite limit at both ends is bounded, so there is a
constant $C$ with $|g_{kl}(\omega)|\le C/(1+\omega^2)$ for every real $\omega$. The comparison function
$|\omega|^{-\beta}/(1+\omega^2)$ is integrable on $\R$ for $-1<\beta<1$, since it behaves like $|\omega|^{-\beta}$ near the origin,
where $-\beta>-1$, and like $|\omega|^{-2-\beta}$ for large $|\omega|$, where $2+\beta>1$. Each integral of $|\omega|^{-\beta}g_{kl}$
therefore converges absolutely.

The noise covariance enters only through $\mathbf{M}_w$, and the scalar integrals $q_{kl}$ depend only on the eigenvalues of
$\mathbf{A}$ and on $\beta$.
The entrywise convergence proved above also establishes convergence of the original covariance integral for diagonalizable
$\mathbf A$. Its value is real because $\mathbf R(-\omega)=\overline{\mathbf R(\omega)}$ and opposite frequencies contribute
complex conjugates. Moreover, $\mathbf R\boldsymbol{\Sigma}_w\mathbf R^\dagger$ is Hermitian positive semidefinite.
Thus $\boldsymbol{\Sigma}$ is real symmetric positive semidefinite; it is positive definite if
$\boldsymbol{\Sigma}_w\succ0$, since $\mathbf R(\omega)$ is invertible for every real $\omega$.

% =====================================================================
\subsection*{Evaluating the scalar integrals in \ref{eq:Sigma_hadamard}}

% =====================================================================

Write $\lambda:=\lambda_k$ and $\mu:=\bar\lambda_l$, both with negative real part. We have
\begin{equation}
    (i\omega-\lambda)(-i\omega-\lambda)=\omega^2+\lambda^2
    \label{eq:lambda_denominator}
\end{equation}
and
\begin{equation}
    (-i\omega-\mu)(i\omega-\mu)=\omega^2+\mu^2.
    \label{eq:mu_denominator}
\end{equation}
Multiplying the numerator and denominator of $g_{kl}$ in \eqref{eq:g_def} by $(-i\omega-\lambda)(i\omega-\mu)$ gives
\begin{equation}
    g_{kl}(\omega)=\frac{(-i\omega-\lambda)(i\omega-\mu)}{(\omega^2+\lambda^2)(\omega^2+\mu^2)}
    =\frac{\omega^2+\lambda\mu+i\omega(\mu-\lambda)}{(\omega^2+\lambda^2)(\omega^2+\mu^2)},
    \label{eq:g_rational}
\end{equation}
The term with $i\omega(\mu-\lambda)$ is odd in $\omega$. Multiplied by $|\omega|^{-\beta}$, it is absolutely integrable, since it is
$O(|\omega|^{1-\beta})$ near $\omega=0$, where the denominator is $\lambda^2\mu^2\neq0$, and $O(|\omega|^{-3-\beta})$ for large
$|\omega|$. Its integral therefore vanishes, leaving
\begin{equation}
    \begin{aligned}
    q_{kl}
    &=\int_{-\infty}^{\infty}|\omega|^{-\beta}g_{kl}(\omega)\,d\omega\\
    &=\int_{-\infty}^{\infty}|\omega|^{-\beta}\,
    \frac{\omega^2+\lambda\mu}{(\omega^2+\lambda^2)(\omega^2+\mu^2)}\,d\omega\\
    &=\frac{1}{\lambda+\mu}\int_{-\infty}^{\infty}|\omega|^{-\beta}
    \Bigl[\frac{\lambda}{\omega^2+\lambda^2}+\frac{\mu}{\omega^2+\mu^2}\Bigr]d\omega.
    \end{aligned}
    \label{eq:q_reduced}
\end{equation}
The sum $\lambda+\mu=\lambda_k+\bar\lambda_l$ has negative real part, so it is nonzero, and
\begin{equation}
    \frac{\omega^2+\lambda\mu}{(\omega^2+\lambda^2)(\omega^2+\mu^2)}
    =\frac{1}{\lambda+\mu}\Bigl[\frac{\lambda}{\omega^2+\lambda^2}+\frac{\mu}{\omega^2+\mu^2}\Bigr],
    \label{eq:pf_general}
\end{equation}
Therefore, the integral in \eqref{eq:q_reduced} can be written equivalently as
\begin{equation}
    q_{kl}=\frac{\lambda I_\beta(\lambda)+\mu I_\beta(\mu)}{\lambda+\mu},
    \qquad
    I_\beta(c):=\int_{\mathbb R}\frac{|\omega|^{-\beta}}{\omega^2+c^2}\,d\omega.
    \label{eq:q_two_integrals}
\end{equation}
For every complex $c$ with
$\operatorname{Re}c<0$, the elementary integral \eqref{A-eq:D_elem} of Appendix~A, extended there to the left half-plane, gives
\begin{equation}
    \int_{-\infty}^{\infty}\frac{|\omega|^{-\beta}}{c^2+\omega^2}\,d\omega
    =\frac{\pi\,(-c)^{-1-\beta}}{\cos(\pi\beta/2)}.
    \label{eq:Ibeta_eval}
\end{equation}
Here the power of $-c$, which lies in the right half-plane, is the principal power, and the integral converges absolutely. The numerators
$\lambda$ and $\mu$ of \eqref{eq:q_two_integrals} lie in the left half-plane, whereas the integral returns a power of
the right half-plane quantity $-c$. Reconciling the two introduces a sign, because writing $c=-(-c)$ turns the leading factor
into a power of the same base $-c$,
\begin{equation}
    c\,(-c)^{-1-\beta}=-(-c)^{1}\,(-c)^{-1-\beta}=-(-c)^{-\beta} ,
    \label{eq:power_sign}
\end{equation}
where the exponents add to $1+(-1-\beta)=-\beta$ because both powers use the same principal branch of the logarithm of $-c$.
Applying this with $c=\lambda$ and $c=\mu$ to the two terms of \eqref{eq:q_two_integrals} therefore gives
\begin{equation}
\begin{aligned}
    q_{kl}&=\int_{-\infty}^{\infty}|\omega|^{-\beta}\,
    \frac{\omega^2+\lambda\mu}{(\omega^2+\lambda^2)(\omega^2+\mu^2)}\,d\omega\\
    &=\frac{1}{\lambda+\mu}\int_{-\infty}^{\infty}|\omega|^{-\beta}
    \Bigl[\frac{\lambda}{\omega^2+\lambda^2}+\frac{\mu}{\omega^2+\mu^2}\Bigr]d\omega\\
    &=\frac{1}{\lambda+\mu}\Bigl[\frac{\pi\,\lambda\,(-\lambda)^{-1-\beta}}{\cos(\pi\beta/2)}+\frac{\pi\,\mu\,(-\mu)^{-1-\beta}}{\cos(\pi\beta/2)}\Bigr]\\
    &=\frac{-\pi\bigl[(-\lambda)^{-\beta}+(-\mu)^{-\beta}\bigr]}{(\lambda+\mu)\cos(\pi\beta/2)}\\
    &=\frac{\pi}{\cos(\pi\beta/2)}\cdot\frac{(-\lambda_k)^{-\beta}+(-\bar\lambda_l)^{-\beta}}{-(\lambda_k+\bar\lambda_l)} .
\end{aligned}
    \label{eq:q_closed}
\end{equation}
Both terms contribute the same factor $-1$, which is why it factors out of the bracket in the first line. The second line
carries that factor into the denominator, since $-1/(\lambda+\mu)=1/\bigl(-(\lambda+\mu)\bigr)$, and restores
$\lambda=\lambda_k$ and $\mu=\bar\lambda_l$. Every base in the result then has positive real part, because
$\operatorname{Re}\lambda_k<0$ and $\operatorname{Re}\lambda_l<0$ give $\operatorname{Re}\bigl(-(\lambda_k+\bar\lambda_l)\bigr)>0$.
Finally, the identity in \eqref{eq:pf_general} does not require $\lambda$ and $\mu$ to be distinct. When $\lambda=\mu$, both sides reduce directly to
\begin{equation}
    \frac{\omega^2+\lambda^2}{(\omega^2+\lambda^2)^2}
    =\frac{1}{2\lambda}\Bigl[\frac{\lambda}{\omega^2+\lambda^2}+\frac{\lambda}{\omega^2+\lambda^2}\Bigr]
    =\frac{1}{\omega^2+\lambda^2}.
    \label{eq:pf_repeated}
\end{equation}
More generally, if $\lambda_k=\lambda_l$ is a repeated eigenvalue, then
$\lambda_k+\bar\lambda_l=2\operatorname{Re}\lambda_k<0$.

% =====================================================================
\subsection*{The covariance formula satisfies a Lyapunov equation}
% =====================================================================

Substituting \eqref{eq:q_closed} into \eqref{eq:Sigma_hadamard} gives, for every diagonalizable $\mathbf{A}\in\mathcal{U}$,
\begin{equation}
    \boldsymbol{\Sigma}(\mathbf{A},\beta)=\frac{\alpha}{2\cos(\pi\beta/2)}\,\mathbf{P}\,
    \Bigl[(M_w)_{kl}\,\frac{(-\lambda_k)^{-\beta}+(-\bar\lambda_l)^{-\beta}}{-(\lambda_k+\bar\lambda_l)}\Bigr]_{k,l=1}^{n}\,\mathbf{P}^{\dagger},
    \label{eq:Sigma_closed_Sw}
\end{equation}
where $[x_{kl}]_{k,l=1}^n$ denotes the $n\times n$ matrix with entries $x_{kl}$. The right-hand side is real, since it equals
$\boldsymbol{\Sigma}(\mathbf{A},\beta)$, although $\mathbf{P}$, $\mathbf{M}_w$ and the $\lambda_k$ may be complex. For the isotropic
spectrum, $\boldsymbol{\Sigma}_w=\mathbf{I}$ gives $\mathbf{M}_w=\mathbf{M}:=\mathbf{P}^{-1}\mathbf{P}^{-\dagger}$, and \eqref{eq:Sigma_closed_Sw} becomes
\begin{equation}
    \boldsymbol{\Sigma}(\mathbf{A},\beta)=\frac{\alpha}{2\cos(\pi\beta/2)}\,\mathbf{P}\,
    \Bigl[M_{kl}\,\frac{(-\lambda_k)^{-\beta}+(-\bar\lambda_l)^{-\beta}}{-(\lambda_k+\bar\lambda_l)}\Bigr]_{k,l=1}^{n}\,\mathbf{P}^{\dagger}.
    \label{eq:Sigma_closed}
\end{equation}

\textbf{The Lyapunov operator $(\mathbf A\boldsymbol\Sigma+\boldsymbol\Sigma\mathbf A^\top)$ in the eigenbasis.}
Because $\mathbf A=\mathbf P\boldsymbol\Lambda\mathbf P^{-1}$ and $\mathbf A$ is real,
$\mathbf A^\top=\mathbf A^\dagger=\mathbf P^{-\dagger}\bar{\boldsymbol\Lambda}\mathbf P^\dagger$.
Consequently,
\begin{equation}
\begin{aligned}
    \mathbf P^{-1}(\mathbf A\boldsymbol\Sigma+\boldsymbol\Sigma\mathbf A^\top)\mathbf P^{-\dagger}
    &=\mathbf P^{-1}\Bigl[(\mathbf P\boldsymbol\Lambda\mathbf P^{-1})\boldsymbol\Sigma
      +\boldsymbol\Sigma(\mathbf P^{-\dagger}\bar{\boldsymbol\Lambda}\mathbf P^\dagger)\Bigr]\mathbf P^{-\dagger}\\
    &=\mathbf P^{-1}\mathbf P\boldsymbol\Lambda\mathbf P^{-1}\boldsymbol\Sigma\mathbf P^{-\dagger}
      +\mathbf P^{-1}\boldsymbol\Sigma\mathbf P^{-\dagger}\bar{\boldsymbol\Lambda}\mathbf P^\dagger\mathbf P^{-\dagger}\\
    &=\boldsymbol\Lambda\mathbf P^{-1}\boldsymbol\Sigma\mathbf P^{-\dagger}
      +\mathbf P^{-1}\boldsymbol\Sigma\mathbf P^{-\dagger}\bar{\boldsymbol\Lambda}\\
    &=\boldsymbol\Lambda\tilde{\boldsymbol\Sigma}
      +\tilde{\boldsymbol\Sigma}\bar{\boldsymbol\Lambda}.
\end{aligned}
    \label{eq:lyap_transform_eigenbasis}
\end{equation}
Multiplying \eqref{eq:Sigma_closed_Sw} by $\mathbf{P}^{-1}$ on the left and by $\mathbf{P}^{-\dagger}$ on the right, and using
$\mathbf{P}^{\dagger}\mathbf{P}^{-\dagger}=(\mathbf{P}^{-1}\mathbf{P})^{\dagger}=\mathbf{I}$, we have
$\tilde{\boldsymbol{\Sigma}}:=\mathbf{P}^{-1}\boldsymbol{\Sigma}\mathbf{P}^{-\dagger}$, whose $(k,l)$ entry is
\begin{equation}
    \tilde\Sigma_{kl}=\eta_\beta\,(M_w)_{kl}
    \frac{(-\lambda_k)^{-\beta}+(-\bar\lambda_l)^{-\beta}}
    {-(\lambda_k+\bar\lambda_l)},
    \label{eq:Sigma_tilde_entries}
\end{equation}
where $\eta_\beta:=\alpha/\bigl(2\cos(\pi\beta/2)\bigr)$.
The $(k,l)$ entry of $\boldsymbol\Lambda\tilde{\boldsymbol\Sigma}
+\tilde{\boldsymbol\Sigma}\bar{\boldsymbol\Lambda}$ is
\begin{equation}
\begin{aligned}
    \bigl(\boldsymbol\Lambda\tilde{\boldsymbol\Sigma}
      +\tilde{\boldsymbol\Sigma}\bar{\boldsymbol\Lambda}\bigr)_{kl}
    &=(\lambda_k+\bar\lambda_l)\tilde\Sigma_{kl}\\
    &=(\lambda_k+\bar\lambda_l)\eta_\beta(M_w)_{kl}
      \frac{(-\lambda_k)^{-\beta}+(-\bar\lambda_l)^{-\beta}}
      {-(\lambda_k+\bar\lambda_l)}\\
    &=-\eta_\beta(M_w)_{kl}
      \bigl[(-\lambda_k)^{-\beta}+(-\bar\lambda_l)^{-\beta}\bigr]\\
    &=-\eta_\beta\Bigl[
      (-\lambda_k)^{-\beta}(M_w)_{kl}
      +(M_w)_{kl}(-\bar\lambda_l)^{-\beta}\Bigr].
\end{aligned}
    \label{eq:lyap_entrywise}
\end{equation}
Since $\operatorname{Re}(\lambda_k+\bar\lambda_l)<0$ by stability, every multiplier
$\lambda_k+\bar\lambda_l$ is nonzero; hence the Lyapunov operator is invertible.
Define
\begin{equation*}
    (-\boldsymbol\Lambda)^{-\beta}
    :=\operatorname{diag}\bigl((-\lambda_1)^{-\beta},\dots,(-\lambda_n)^{-\beta}\bigr),
    \qquad
    (-\bar{\boldsymbol\Lambda})^{-\beta}
    :=\operatorname{diag}\bigl((-\bar\lambda_1)^{-\beta},\dots,(-\bar\lambda_n)^{-\beta}\bigr).
\end{equation*}
Then
\begin{equation}
\begin{aligned}
    (-\lambda_k)^{-\beta}(M_w)_{kl}
    &=\bigl[(-\boldsymbol\Lambda)^{-\beta}\bigr]_{kk}(M_w)_{kl}\\
    &=\bigl[(-\boldsymbol\Lambda)^{-\beta}\mathbf M_w\bigr]_{kl},
\end{aligned}
    \label{eq:left_power_entry}
\end{equation}
and
\begin{equation}
\begin{aligned}
    (M_w)_{kl}(-\bar\lambda_l)^{-\beta}
    &=(M_w)_{kl}\bigl[(-\bar{\boldsymbol\Lambda})^{-\beta}\bigr]_{ll}\\
    &=\bigl[\mathbf M_w(-\bar{\boldsymbol\Lambda})^{-\beta}\bigr]_{kl}.
\end{aligned}
    \label{eq:right_power_entry}
\end{equation}
Therefore, the entrywise identities in \eqref{eq:lyap_entrywise} can be assembled into
\begin{equation}
    \boldsymbol\Lambda\tilde{\boldsymbol\Sigma}+\tilde{\boldsymbol\Sigma}\bar{\boldsymbol\Lambda}
    =-\eta_\beta\Bigl[(-\boldsymbol\Lambda)^{-\beta}\mathbf M_w
    +\mathbf M_w(-\bar{\boldsymbol\Lambda})^{-\beta}\Bigr].
    \label{eq:lyap_eigen_coordinates}
\end{equation}

\textbf{The Lyapunov operator $(\mathbf A\boldsymbol\Sigma+\boldsymbol\Sigma\mathbf A^\top)$ in the original coordinates.}
Multiplying both sides of \eqref{eq:lyap_eigen_coordinates} on the left by $\mathbf{P}$ and on the right by $\mathbf{P}^{\dagger}$.
The left-hand side becomes
\begin{equation}
\begin{aligned}
    \mathbf{P}\bigl(\boldsymbol\Lambda\tilde{\boldsymbol\Sigma}
      +\tilde{\boldsymbol\Sigma}\bar{\boldsymbol\Lambda}\bigr)\mathbf{P}^{\dagger}
    &=\mathbf{P}\Bigl[\mathbf{P}^{-1}
      (\mathbf{A}\boldsymbol\Sigma+\boldsymbol\Sigma\mathbf{A}^{\top})
      \mathbf{P}^{-\dagger}\Bigr]\mathbf{P}^{\dagger}\\
    &=\mathbf{A}\boldsymbol\Sigma+\boldsymbol\Sigma\mathbf{A}^{\top}.
\end{aligned}
    \label{eq:lyap_return_left}
\end{equation}
The right-hand side becomes
\begin{equation}
\begin{aligned}
    &-\eta_\beta\,\mathbf{P}\Bigl[
      (-\boldsymbol{\Lambda})^{-\beta}\mathbf{M}_w
      +\mathbf{M}_w(-\bar{\boldsymbol{\Lambda}})^{-\beta}
      \Bigr]\mathbf{P}^{\dagger}\\
    &=-\eta_\beta\,\mathbf{P}\Bigl[
      (-\boldsymbol{\Lambda})^{-\beta}(\mathbf{P}^{-1}\boldsymbol{\Sigma}_w\mathbf{P}^{-\dagger})
      +(\mathbf{P}^{-1}\boldsymbol{\Sigma}_w\mathbf{P}^{-\dagger})(-\bar{\boldsymbol{\Lambda}})^{-\beta}
      \Bigr]\mathbf{P}^{\dagger}\\
    &=-\eta_\beta\Bigl[
      \bigl[\mathbf{P}(-\boldsymbol{\Lambda})^{-\beta}\mathbf{P}^{-1}\bigr]\boldsymbol{\Sigma}_w
      \bigl[\mathbf{P}^{-\dagger}\mathbf{P}^{\dagger}\bigr]
      +\bigl[\mathbf{P}\mathbf{P}^{-1}\bigr]\boldsymbol{\Sigma}_w
      \bigl[\mathbf{P}^{-\dagger}(-\bar{\boldsymbol{\Lambda}})^{-\beta}\mathbf{P}^{\dagger}\bigr]
      \Bigr]\\
    &=-\eta_\beta\Bigl[
      (-\mathbf{A})^{-\beta}\boldsymbol{\Sigma}_w
      +\boldsymbol{\Sigma}_w\bigl[(-\mathbf{A})^{-\beta}\bigr]^{\dagger}
      \Bigr].
\end{aligned}
    \label{eq:lyap_return_right}
\end{equation}
where $(-\mathbf{A})^{-\beta}:=\mathbf{P}(-\boldsymbol{\Lambda})^{-\beta}\mathbf{P}^{-1}$ is the principal power of $-\mathbf{A}$, and the
last equality uses
$\bigl[\mathbf{P}(-\boldsymbol{\Lambda})^{-\beta}\mathbf{P}^{-1}\bigr]^{\dagger}=\mathbf{P}^{-\dagger}(-\bar{\boldsymbol{\Lambda}})^{-\beta}\mathbf{P}^{\dagger}$.
The principal power does not depend on the choice of eigenvectors, since it multiplies every eigenvector of $\mathbf{A}$ with
eigenvalue $\lambda$ by $(-\lambda)^{-\beta}$, and the eigenvectors span $\mathbb{C}^n$. It is also real. Since $\mathbf{A}$ is real,
its eigenvalues and eigenvectors occur in conjugate pairs, and $\operatorname{Re}(-\lambda_k)>0$ keeps every $-\lambda_k$ off the cut
$(-\infty,0]$ of the principal branch, where $\overline{z^{-\beta}}=\bar z^{-\beta}$. Conjugating $(-\mathbf{A})^{-\beta}$ therefore
permutes its eigenvalues and eigenvectors among themselves and returns the same matrix, so its conjugate transpose is its
transpose. With $\eta_\beta=\alpha/\bigl(2\cos(\pi\beta/2)\bigr)$, \eqref{eq:lyap_eigen_coordinates} becomes
\begin{equation}
\boxed{\mathbf{A}\boldsymbol{\Sigma}+\boldsymbol{\Sigma}\mathbf{A}^\top
    +\eta_\beta\Bigl[(-\mathbf{A})^{-\beta}\boldsymbol{\Sigma}_w
    +\boldsymbol{\Sigma}_w\bigl[(-\mathbf{A})^{-\beta}\bigr]^{\top}\Bigr]=\mathbf{0}}
    \label{eq:lyap_colored_Sw}
\end{equation}
and, for the isotropic spectrum $\boldsymbol{\Sigma}_w=\mathbf{I}$,
\begin{equation}
    \mathbf{A}\boldsymbol{\Sigma}+\boldsymbol{\Sigma}\mathbf{A}^\top
    +\eta_\beta\Bigl[(-\mathbf{A})^{-\beta}+\bigl[(-\mathbf{A})^{-\beta}\bigr]^{\top}\Bigr]=\mathbf{0} .
    \label{eq:lyap_colored}
\end{equation}
Equivalently, the general colored-noise equation \eqref{eq:lyap_colored_Sw} has the standard continuous-time Lyapunov form
\begin{equation}
\mathbf A\boldsymbol\Sigma+\boldsymbol\Sigma\mathbf A^\top+
\mathbf Q_\beta(\mathbf A,\boldsymbol\Sigma_w)=\mathbf0,
\label{eq:lyap_standard_form}
\end{equation}
where
\begin{equation}
\mathbf Q_\beta(\mathbf A,\boldsymbol\Sigma_w)
=\eta_\beta\!\left[
(-\mathbf A)^{-\beta}\boldsymbol\Sigma_w+
\boldsymbol\Sigma_w\bigl[(-\mathbf A)^{-\beta}\bigr]^\top
\right].
\label{eq:Q_beta}
\end{equation}
At $\beta=0$ (\textbf{white noise}), $\eta_0=\alpha/\bigl(2\cos 0\bigr)=\alpha/2$ and
$(-\mathbf{A})^{0}=\mathbf{I}$, so \eqref{eq:lyap_colored} becomes
\begin{equation}
    \mathbf{A}\boldsymbol{\Sigma}+\boldsymbol{\Sigma}\mathbf{A}^{\top}=-\alpha\mathbf{I}.
    \label{eq:lyap_white_noise}
\end{equation}

\textbf{Vectorized solution of the common-exponent Lyapunov equation.}
For a matrix $\mathbf X$, let $\vc(\mathbf X)$ denote column-stacking vectorization, and let $\otimes$ denote the Kronecker product. The identity
\begin{equation}
\vc(\mathbf U\mathbf X\mathbf V)=(\mathbf V^\top\otimes\mathbf U)\vc(\mathbf X)
\label{eq:vec_product_identity}
\end{equation}
holds for conformable real matrices. Set
\begin{equation}
\mathbf B_\beta:=(-\mathbf A)^{-\beta},
\qquad
\mathbf K(\mathbf A):=\mathbf I\otimes\mathbf A+\mathbf A\otimes\mathbf I.
\label{eq:kronecker_common_definitions}
\end{equation}
Then the two terms containing the unknown covariance in \eqref{eq:lyap_colored_Sw} satisfy
\begin{equation}
\begin{aligned}
\vc(\mathbf A\boldsymbol\Sigma+\boldsymbol\Sigma\mathbf A^\top)
&=\vc(\mathbf A\boldsymbol\Sigma)+\vc(\boldsymbol\Sigma\mathbf A^\top)\\
&=(\mathbf I\otimes\mathbf A)\vc(\boldsymbol\Sigma)
 +(\mathbf A\otimes\mathbf I)\vc(\boldsymbol\Sigma)\\
&=\mathbf K(\mathbf A)\vc(\boldsymbol\Sigma).
\end{aligned}
\label{eq:vec_lyapunov_operator}
\end{equation}
whereas its forcing terms satisfy
\begin{equation}
\begin{aligned}
\vc(\mathbf B_\beta\boldsymbol\Sigma_w+\boldsymbol\Sigma_w\mathbf B_\beta^\top)
&=\vc(\mathbf B_\beta\boldsymbol\Sigma_w)
 +\vc(\boldsymbol\Sigma_w\mathbf B_\beta^\top)\\
&=(\mathbf I\otimes\mathbf B_\beta)\vc(\boldsymbol\Sigma_w)
 +(\mathbf B_\beta\otimes\mathbf I)\vc(\boldsymbol\Sigma_w)\\
&=\bigl(\mathbf I\otimes\mathbf B_\beta
 +\mathbf B_\beta\otimes\mathbf I\bigr)\vc(\boldsymbol\Sigma_w).
\end{aligned}
\label{eq:vec_common_forcing}
\end{equation}
Consequently, \eqref{eq:lyap_colored_Sw} is equivalent to
\begin{equation}
\mathbf K(\mathbf A)\vc(\boldsymbol\Sigma)
=-\eta_\beta\bigl(\mathbf I\otimes\mathbf B_\beta
+\mathbf B_\beta\otimes\mathbf I\bigr)\vc(\boldsymbol\Sigma_w).
\label{eq:vec_common_linear_system}
\end{equation}
The eigenvalues of $\mathbf K(\mathbf A)$ are $\lambda_r+\lambda_s$ for $r,s=1,\ldots,n$.
Their real parts are strictly negative because $\mathbf A$ is Hurwitz stable, so $\mathbf K(\mathbf A)$ is invertible. Therefore,
\begin{equation}
\boxed{\begin{aligned}
\vc(\boldsymbol\Sigma)
&=-\mathbf K(\mathbf A)^{-1}
\left\{\eta_\beta\bigl[\mathbf I\otimes(-\mathbf A)^{-\beta}
+(-\mathbf A)^{-\beta}\otimes\mathbf I\bigr]\right\}
\vc(\boldsymbol\Sigma_w)\\
&=-\bigl(\mathbf I\otimes\mathbf A
+\mathbf A\otimes\mathbf I\bigr)^{-1}
\left\{\eta_\beta\bigl[\mathbf I\otimes(-\mathbf A)^{-\beta}
+(-\mathbf A)^{-\beta}\otimes\mathbf I\bigr]\right\}
\vc(\boldsymbol\Sigma_w).
\end{aligned}}
\label{eq:vec_common_solution}
\end{equation}
\subsection*{Interpreting the $(-\mathbf A)^{-\beta}\boldsymbol{\Sigma}_w$ term in the Lyapunov equation}

Here we show that the forcing term in the Lyapunov equation \eqref{eq:lyap_colored_Sw} is a sum of the state--noise cross-covariance and its transpose. Define the formal zero-lag cross-covariance by the symmetric-limit integral
\begin{equation}
    \boldsymbol{\Sigma}_{xw}
    :=\operatorname{p.v.}\frac{1}{2\pi}\int_{\mathbb R}
    \mathbf R(\omega;\mathbf A)\,\mathbf S_w(\omega)\,d\omega .
    \label{eq:Sigma_xw_definition}
\end{equation}
The principal-value interpretation is needed because $\boldsymbol w$ is generalized noise. For $\beta>0$ the integral is absolutely convergent at high frequency; for $\beta\le0$, its leading odd term is canceled by the symmetric limit. We first evaluate this cross-covariance. Since $\mathbf S_w(\omega)$ is even in $\omega$,
\begin{equation}
\boldsymbol{\Sigma}_{xw}
=\frac{\alpha}{2\pi}\int_{0}^{\infty}\omega^{-\beta}
\bigl[\mathbf R(\omega;\mathbf A)+\mathbf R(-\omega;\mathbf A)\bigr]\,d\omega\,
\boldsymbol{\Sigma}_w .
\label{eq:Sigma_xw_even}
\end{equation}
The two resolvents can be combined without assuming that $\mathbf A$ is symmetric:
\begin{equation}
\begin{aligned}
\mathbf R(\omega;\mathbf A)+\mathbf R(-\omega;\mathbf A)
&=\bigl[(-i\omega\mathbf I-\mathbf A)+(i\omega\mathbf I-\mathbf A)\bigr]
  (\omega^2\mathbf I+\mathbf A^2)^{-1}\\
&=-2\mathbf A(\omega^2\mathbf I+\mathbf A^2)^{-1}.
\end{aligned}
\label{eq:resolvent_pair_sum}
\end{equation}
Here we used the commuting factorization
\begin{equation}
(i\omega\mathbf I-\mathbf A)(-i\omega\mathbf I-\mathbf A)
=\omega^2\mathbf I+\mathbf A^2 .
\label{eq:resolvent_product}
\end{equation}
The factor is invertible for real $\omega$ because its eigenvalues
$(i\omega-\lambda_k)(-i\omega-\lambda_k)$ are nonzero when
$\operatorname{Re}\lambda_k<0$. Hence
\begin{equation}
\boldsymbol{\Sigma}_{xw}
=\frac{\alpha}{\pi}\int_{0}^{\infty}\omega^{-\beta}
(-\mathbf A)(\omega^2\mathbf I+\mathbf A^2)^{-1}\,d\omega\,
\boldsymbol{\Sigma}_w .
\label{eq:Sigma_xw_reduced}
\end{equation}

Now use $\mathbf A=\mathbf P\boldsymbol\Lambda\mathbf P^{-1}$. First,
\begin{equation}
\mathbf A^2
=(\mathbf P\boldsymbol\Lambda\mathbf P^{-1})(\mathbf P\boldsymbol\Lambda\mathbf P^{-1})
=\mathbf P\boldsymbol\Lambda^2\mathbf P^{-1},
\label{eq:Sigma_xw_A_squared}
\end{equation}
where $\mathbf P^{-1}\mathbf P=\mathbf I$. Therefore,
\begin{equation}
\omega^2\mathbf I+\mathbf A^2
=\mathbf P(\omega^2\mathbf I+\boldsymbol\Lambda^2)\mathbf P^{-1},
\label{eq:Sigma_xw_shifted}
\end{equation}
and inversion reverses the similarity factors:
\begin{equation}
(\omega^2\mathbf I+\mathbf A^2)^{-1}
=\mathbf P(\omega^2\mathbf I+\boldsymbol\Lambda^2)^{-1}\mathbf P^{-1}.
\label{eq:Sigma_xw_inverse}
\end{equation}
Since $-\mathbf A=\mathbf P(-\boldsymbol\Lambda)\mathbf P^{-1}$, multiplying the last two identities gives
\begin{equation}
\begin{aligned}
(-\mathbf A)(\omega^2\mathbf I+\mathbf A^2)^{-1}
&=\mathbf P(-\boldsymbol\Lambda)(\omega^2\mathbf I+\boldsymbol\Lambda^2)^{-1}\mathbf P^{-1}\\
&=\mathbf P\operatorname{diag}\!\left(
\frac{-\lambda_1}{\omega^2+\lambda_1^2},\ldots,
\frac{-\lambda_n}{\omega^2+\lambda_n^2}\right)\mathbf P^{-1},
\end{aligned}
\label{eq:Sigma_xw_diagonal}
\end{equation}
because $\boldsymbol\Lambda$ is diagonal.
For each eigenvalue, the scalar integral established earlier gives
\begin{equation}
\int_{0}^{\infty}\omega^{-\beta}
\frac{-\lambda_k}{\omega^2+\lambda_k^2}\,d\omega
=\frac{\pi}{2\cos(\pi\beta/2)}(-\lambda_k)^{-\beta}.
\label{eq:Sigma_xw_scalar}
\end{equation}
Therefore, with $\eta_\beta=\alpha/[2\cos(\pi\beta/2)]$,
\begin{equation}
    \boldsymbol{\Sigma}_{xw}
    =\eta_\beta(-\mathbf A)^{-\beta}\boldsymbol{\Sigma}_w .
    \label{eq:Sigma_xw_closed}
\end{equation}

With this identification, the forcing term in the covariance equation is
\begin{equation}
    \mathbf A\boldsymbol{\Sigma}+\boldsymbol{\Sigma}\mathbf A^\top
    +\boldsymbol{\Sigma}_{xw}+\boldsymbol{\Sigma}_{xw}^{\top}
    =\mathbf 0 .
    \label{eq:lyap_cross_covariance}
\end{equation}

% =====================================================================
\subsection*{Extending the results to $n \times n$ networks with region-specific noise exponents}
% =====================================================================

Let region $j$ have exponent $\beta_j\in(-1,1)$, and retain the real symmetric positive-semidefinite spatial covariance factor
$\boldsymbol\Sigma_w$ from the setting above. It may have nonzero off-diagonal entries, representing spatially correlated
neuronal inputs. Define
\begin{equation}
\mathbf D(\omega):=\operatorname{diag}\bigl(
|\omega|^{-\beta_1/2},\ldots,|\omega|^{-\beta_n/2}\bigr),
\qquad
\mathbf S_w(\omega)
=\alpha\,\mathbf D(\omega)\boldsymbol{\Sigma}_w\mathbf D(\omega).
\label{eq:regional_noise_spectrum}
\end{equation}
For every nonzero frequency, $\mathbf D(\omega)$ is real and diagonal, but $\mathbf S_w(\omega)$ need not be diagonal.
It is positive semidefinite because, for every $\mathbf v\in\mathbb C^n$,
$\mathbf v^\dagger\mathbf S_w(\omega)\mathbf v
=\alpha[\mathbf D(\omega)\mathbf v]^\dagger\boldsymbol\Sigma_w[\mathbf D(\omega)\mathbf v]\geq0$.
The independent-input model of \href{https://doi.org/10.1162/netn_a_00348}{\citet{Novelli2024}} is the special case in which $\boldsymbol\Sigma_w$ is diagonal.

Write $\sigma_{jk}:=(\boldsymbol\Sigma_w)_{jk}=\sigma_{kj}$, and, for $1\leq j\leq k\leq n$, define
\begin{equation}
\begin{gathered}
\beta_{jk}:=\frac{\beta_j+\beta_k}{2},\qquad
\mathbf Q_{jk}:=
\begin{cases}
\sigma_{jj}\mathbf e_j\mathbf e_j^\top,&j=k,\\
\sigma_{jk}(\mathbf e_j\mathbf e_k^\top+\mathbf e_k\mathbf e_j^\top),&j<k.
\end{cases}\\
\boldsymbol\Sigma_w=\sum_{j\leq k}\mathbf Q_{jk}.
\end{gathered}
\label{eq:regional_pair_definitions}
\end{equation}
Indeed, $\mathbf D(\omega)\mathbf e_j=|\omega|^{-\beta_j/2}\mathbf e_j$. Thus each
$(j,k)$ entry of \eqref{eq:regional_noise_spectrum} has exponent
$\beta_j/2+\beta_k/2=\beta_{jk}$ and equals
$\alpha\sigma_{jk}|\omega|^{-\beta_{jk}}$. Grouping symmetric entries gives
\begin{equation}
\mathbf S_w(\omega)
=\alpha\sum_{j\leq k}|\omega|^{-\beta_{jk}}\mathbf Q_{jk}.
\label{eq:regional_noise_expansion}
\end{equation}
Each $\mathbf Q_{jk}$ is real symmetric, but an off-diagonal $\mathbf Q_{jk}$ need not be positive semidefinite.
This does not affect the algebra: every real symmetric matrix is a difference of two positive-semidefinite matrices,
and the covariance integral, its closed form \eqref{eq:Sigma_closed_Sw}, and the Lyapunov equation
\eqref{eq:lyap_colored_Sw} are linear in the spatial factor for a fixed exponent. We apply them to each
$\mathbf Q_{jk}$ by this linear extension. Positive semidefiniteness applies to the complete spectrum and
the complete covariance, not necessarily to individual pair terms.

For each pair, set
\begin{equation}
\mathbf M^{(jk)}:=\mathbf P^{-1}\mathbf Q_{jk}\mathbf P^{-\dagger},
\qquad
\boldsymbol{\Sigma}^{(jk)}
:=\frac{\alpha}{2\pi}\int_{\mathbb R}|\omega|^{-\beta_{jk}}
\mathbf R(\omega;\mathbf A)\mathbf Q_{jk}\mathbf R(\omega;\mathbf A)^\dagger\,d\omega .
\label{eq:regional_covariance_parts}
\end{equation}
Each $\beta_{jk}$ lies in $(-1,1)$. Hurwitz stability bounds the resolvent near zero frequency, giving
an integrable $|\omega|^{-\beta_{jk}}$ factor there; at high frequency, the two resolvents give
$O(|\omega|^{-2-\beta_{jk}})$. Thus every pair integral converges absolutely, even if its matrix
$\mathbf Q_{jk}$ is indefinite. By \eqref{eq:regional_noise_expansion},
$\boldsymbol{\Sigma}=\sum_{j\leq k}\boldsymbol{\Sigma}^{(jk)}$. Applying the closed form
\eqref{eq:Sigma_closed_Sw} term by term gives
\begin{equation}
\begin{aligned}
\boldsymbol{\Sigma}(\mathbf A;\beta_1,\ldots,\beta_n)
={}&\sum_{j\leq k}\eta_{\beta_{jk}}\,\mathbf P
\left[
(\mathbf M^{(jk)})_{ab}
\frac{(-\lambda_a)^{-\beta_{jk}}+(-\bar\lambda_b)^{-\beta_{jk}}}
{-(\lambda_a+\bar\lambda_b)}
\right]_{a,b=1}^{n}
\mathbf P^\dagger,
\end{aligned}
\label{eq:Sigma_closed_regional}
\end{equation}
where $\eta_{\beta_{jk}}:=\alpha/[2\cos(\pi\beta_{jk}/2)]$ and every scalar and matrix power uses
the principal branch specified above.

\textbf{The Lyapunov equation for regional exponents and spatial correlations.}
By the same linearity, \eqref{eq:lyap_colored_Sw} applied to each pair contribution gives
\begin{equation}
\mathbf A\boldsymbol{\Sigma}^{(jk)}+\boldsymbol{\Sigma}^{(jk)}\mathbf A^\top
+\eta_{\beta_{jk}}\left[
(-\mathbf A)^{-\beta_{jk}}\mathbf Q_{jk}
+\mathbf Q_{jk}[(-\mathbf A)^{-\beta_{jk}}]^\top
\right]=\mathbf0.
\label{eq:regional_lyapunov_part}
\end{equation}
Summing over all pairs yields
\begin{equation}
\boxed{
\begin{aligned}
\mathbf A\boldsymbol{\Sigma}+\boldsymbol{\Sigma}\mathbf A^\top
+\sum_{j\leq k}\eta_{\beta_{jk}}\left[
(-\mathbf A)^{-\beta_{jk}}\mathbf Q_{jk}
+\mathbf Q_{jk}[(-\mathbf A)^{-\beta_{jk}}]^\top
\right]=\mathbf0 .
\end{aligned}}
\label{eq:lyap_regional}
\end{equation}
\textbf{Vectorized solution of the regional Lyapunov equation.}
The same Kronecker representation applies to \eqref{eq:lyap_regional}. Let
\begin{equation}
\mathbf B_{jk}:=(-\mathbf A)^{-\beta_{jk}},
\qquad
\mathbf E_{jk}:=\mathbf e_j\mathbf e_k^\top,
\label{eq:regional_vector_definitions}
\end{equation}
and define the pair-selection matrices
\begin{equation}
\mathbf C_{jk}:=
\begin{cases}
\vc(\mathbf E_{jj})\vc(\mathbf E_{jj})^\top,&j=k,\\
\vc(\mathbf E_{jk}+\mathbf E_{kj})\vc(\mathbf E_{jk})^\top,&j<k.
\end{cases}
\label{eq:regional_pair_selector}
\end{equation}
Because $\boldsymbol\Sigma_w$ is symmetric, its $(j,k)$ entry is selected by
$\vc(\mathbf E_{jk})^\top\vc(\boldsymbol\Sigma_w)=\sigma_{jk}$. Hence the pair matrices in
\eqref{eq:regional_pair_definitions} obey
\begin{equation}
\vc(\mathbf Q_{jk})=\mathbf C_{jk}\vc(\boldsymbol\Sigma_w),
\qquad j\leq k.
\label{eq:vec_regional_pair}
\end{equation}
Vectorizing \eqref{eq:lyap_regional} gives
\begin{equation}
\begin{aligned}
\mathbf 0
={}&\mathbf K(\mathbf A)\vc(\boldsymbol\Sigma)
+\sum_{j\leq k}\eta_{\beta_{jk}}
\bigl(\mathbf I\otimes\mathbf B_{jk}
+\mathbf B_{jk}\otimes\mathbf I\bigr)\vc(\mathbf Q_{jk})\\
={}&\mathbf K(\mathbf A)\vc(\boldsymbol\Sigma)
+\sum_{j\leq k}\eta_{\beta_{jk}}
\bigl(\mathbf I\otimes\mathbf B_{jk}
+\mathbf B_{jk}\otimes\mathbf I\bigr)\mathbf C_{jk}
\vc(\boldsymbol\Sigma_w).
\end{aligned}
\label{eq:vec_regional_linear_system}
\end{equation}
Since $\mathbf K(\mathbf A)$ is invertible, the covariance is
\begin{equation}
\boxed{\begin{aligned}
\vc(\boldsymbol\Sigma)
&=-\mathbf K(\mathbf A)^{-1}
\left\{
\sum_{j\leq k}\eta_{\beta_{jk}}
\bigl[\mathbf I\otimes(-\mathbf A)^{-\beta_{jk}}
+(-\mathbf A)^{-\beta_{jk}}\otimes\mathbf I\bigr]\mathbf C_{jk}
\right\}
\vc(\boldsymbol\Sigma_w)\\
&=-\bigl(\mathbf I\otimes\mathbf A
+\mathbf A\otimes\mathbf I\bigr)^{-1}
\left\{
\sum_{j\leq k}\eta_{\beta_{jk}}
\bigl[\mathbf I\otimes(-\mathbf A)^{-\beta_{jk}}
+(-\mathbf A)^{-\beta_{jk}}\otimes\mathbf I\bigr]\mathbf C_{jk}
\right\}
\vc(\boldsymbol\Sigma_w).
\end{aligned}}
\label{eq:vec_regional_solution}
\end{equation}
This equation holds for every diagonalizable real Hurwitz-stable $\mathbf A$ in the scope of the eigenbasis derivation.
If $\boldsymbol\Sigma_w$ is diagonal, the off-diagonal pair terms vanish and the independent-input regional
formulas are recovered. If every $\beta_j=\beta$, then $\beta_{jk}=\beta$ for all pairs and
$\sum_{j\leq k}\mathbf Q_{jk}=\boldsymbol{\Sigma}_w$. Thus
\eqref{eq:Sigma_closed_regional} and \eqref{eq:lyap_regional} reduce respectively to
\eqref{eq:Sigma_closed_Sw} and \eqref{eq:lyap_colored_Sw}. If every $\beta_j=0$, then
\eqref{eq:lyap_regional} reduces to the white-noise Lyapunov equation
$\mathbf A\boldsymbol{\Sigma}+\boldsymbol{\Sigma}\mathbf A^\top=-\alpha\boldsymbol{\Sigma}_w$;
for $\boldsymbol\Sigma_w=\mathbf I$, this is \eqref{eq:lyap_white_noise}.

\end{document}
